\phantomsection
\section*{Глава 1. Анализ предметной области}
\addcontentsline{toc}{section}{Глава 1. Анализ предметной области}
\stepcounter{section}

\subsection{Обзор существующих систем и программных продуктов}

Библиотека как объект автоматизации включает несколько взаимосвязанных процессов:
учет книжного фонда, регистрацию читателей, выдачу и возврат книг, контроль
задолженностей, поиск изданий, а также формирование отчетов по движению фонда.
При ручном ведении учета, например в бумажных журналах или отдельных электронных
таблицах, возникают типовые проблемы: дублирование записей, сложность поиска
нужного издания, отсутствие оперативной статистики, риск потери истории выдач,
а также невозможность разграничить доступ между сотрудником библиотеки и обычным
читателем.

В рамках данной работы под небольшой библиотекой понимается библиотека с
ограниченным объемом книжного фонда, небольшим количеством читателей и одним или
несколькими сотрудниками, выполняющими функции библиотекаря. Для такой библиотеки
основными задачами являются учет книг, регистрация читателей, оформление выдачи и
возврата книг, поиск изданий и получение простой статистики. При этом сложные
механизмы профессиональной каталогизации, комплектования фонда, интеграции с
внешними библиотечными системами и распределения работы между несколькими
специализированными отделами могут быть избыточными.

Для решения задач автоматизации библиотечного учета существуют
автоматизированные библиотечно-информационные системы. Они позволяют вести
электронные каталоги, учитывать читателей, фиксировать выдачу и возврат книг,
а также формировать отчеты. Рассмотрим несколько характерных решений: Koha,
ИРБИС64 и 1С:Библиотека.

Koha является свободной и открытой библиотечной системой, предназначенной для
автоматизации основных библиотечных процессов. Система включает средства ведения
каталога, обслуживания читателей, учета выдач и формирования отчетов
\cite{koha_community}. Данное решение является функционально развитым и может
использоваться различными типами библиотек. При этом внедрение Koha требует
настройки серверной инфраструктуры, а также последующего администрирования и
сопровождения системы.

САБ ИРБИС64 представляет собой российское решение для автоматизации
библиотечно-информационных технологических процессов. В состав системы входят
серверная часть и автоматизированные рабочие места, включая каталогизацию,
комплектование, книговыдачу и другие модули \cite{irbis64_elnit}. Данная система
ориентирована на профессиональную библиотечную деятельность и поддерживает широкий
набор специализированных процессов. Для небольшой библиотеки такая функциональность
может быть избыточной, так как требует освоения отдельных модулей и настройки
системы под конкретные процессы.

Решение 1С:Библиотека предназначено для автоматизации деятельности библиотек
разных типов и назначения \cite{onec_library}. Его особенностью является работа
в экосистеме 1С и возможность использования типовых учетных механизмов. Такое
решение удобно для организаций, уже использующих платформу 1С, однако требует
наличия соответствующей инфраструктуры и сопровождения. Кроме того, зависимость
от платформы 1С снижает гибкость при разработке индивидуального веб-интерфейса.

Сравнение рассмотренных решений по основным критериям приведено в
таблице~\ref{tab:competitors_library}.

{
\coursetablesetup
\setlength{\LTleft}{0pt}
\setlength{\LTright}{0pt}

\begin{longtable}{|
p{\dimexpr0.12\textwidth-2\tabcolsep-\arrayrulewidth\relax}|
p{\dimexpr0.22\textwidth-2\tabcolsep-\arrayrulewidth\relax}|
p{\dimexpr0.20\textwidth-2\tabcolsep-\arrayrulewidth\relax}|
p{\dimexpr0.21\textwidth-2\tabcolsep-\arrayrulewidth\relax}|
p{\dimexpr0.25\textwidth-2\tabcolsep-\arrayrulewidth\relax}|}
    \caption{Сравнение существующих систем автоматизации библиотек}
    \label{tab:competitors_library} \\

    \hline
    \textbf{Система} &
    \textbf{Назначение} &
    \textbf{Сложность установки} &
    \textbf{Сложность сопровождения} &
    \textbf{Применимость для небольшой библиотеки} \\ \hline
    \endfirsthead

    \tabcont{\ref{tab:competitors_library}}
    \textbf{Система} &
    \textbf{Назначение} &
    \textbf{Сложность установки} &
    \textbf{Сложность сопровождения} &
    \textbf{Применимость для небольшой библиотеки} \\ \hline
    \endhead

    \hline
    \endfoot

    Koha &
    Открытая АБИС для комплексной автоматизации библиотек &
    Средняя или высокая: требуется серверное окружение и настройка &
    Средняя или высокая: требуется администрирование, обновление и настройка &
    Подходит, но может быть избыточна для базового учета книг и выдач \\ \hline

    ИРБИС64 &
    Российская профессиональная АБИС для библиотечно-информационных процессов &
    Средняя или высокая: используется серверная часть и отдельные рабочие места &
    Высокая: требуется настройка модулей и сопровождение специализированных процессов &
    Подходит для профессиональных библиотек, но избыточна для базовых процессов учета \\ \hline

    1С:Биб\-лио\-тека &
    Учетное решение для автоматизации библиотек на платформе 1С &
    Средняя: требуется платформа 1С и настройка конфигурации &
    Средняя или высокая: требуется сопровождение платформы и обновление конфигурации &
    Подходит при наличии инфраструктуры 1С, но менее удобна как самостоятельное веб-решение \\ \hline
\end{longtable}
}

Проведенный анализ показывает, что существующие системы в основном ориентированы
на комплексную автоматизацию библиотек. Они содержат развитые средства
каталогизации, поддержки библиотечных стандартов, работы с электронными ресурсами
и формирования отчетности. Такие возможности являются полезными для крупных
библиотек, однако для небольшой библиотеки они могут быть избыточными и приводить
к усложнению внедрения, настройки и дальнейшего сопровождения системы.

В связи с этим целесообразна разработка более простого веб-приложения,
ориентированного на базовые процессы библиотечного учета: ведение списка книг,
учет читателей, регистрацию выдачи и возврата книг, поиск и фильтрацию записей,
а также получение простой статистики. Отличие разрабатываемого решения заключается
в упрощенной структуре, ориентации на наиболее востребованные функции и отсутствии
избыточных модулей, необходимых только для крупных библиотек или профессиональных
библиотечно-информационных систем.

\subsection{Анализ программных инструментов разработки веб-приложений}

Для разработки веб-приложения выбран стек Python, Django, SQLite, HTML, CSS и
JavaScript. Такой набор технологий позволяет реализовать серверную часть,
пользовательский интерфейс, хранение данных, работу с файлами и разграничение
доступа в рамках единой архитектуры.

Python выбран в качестве основного языка программирования серверной части.
К его преимуществам относятся читаемый синтаксис, большое количество библиотек,
активная экосистема и широкое применение в веб-разработке. В качестве альтернативы
мог бы использоваться JavaScript на серверной стороне, однако в таком случае
потребовалась бы отдельная настройка серверной платформы и дополнительных
библиотек для реализации авторизации, работы с базой данных и шаблонами. Язык C++
также может применяться для серверной разработки, но для создания типового
веб-приложения с формами, базой данных и пользовательскими ролями он требует
больше низкоуровневой настройки и менее удобен для быстрой реализации
прикладной бизнес-логики. Поэтому для данного приложения выбран Python как более
удобный язык для разработки серверной логики учета книг, читателей, выдач и
возвратов.

В качестве веб-фреймворка выбран Django. Django является высокоуровневым
Python-фреймворком, ориентированным на быструю разработку веб-приложений и
чистую структуру проекта \cite{django_main}. Одним из важных преимуществ Django
является наличие встроенных механизмов, которые часто требуются в информационных
системах: работа с пользователями, сессиями, правами доступа, формами, базой
данных и административным интерфейсом.

Альтернативой Django является Flask. Flask является микрофреймворком для Python,
который предоставляет базовые средства маршрутизации и обработки запросов
\cite{flask_docs}. Flask удобен для небольших приложений и дает разработчику
больше свободы в выборе компонентов. Однако для реализации полноценной системы
учета с авторизацией, ролями, CRUD-интерфейсом и формами потребовалось бы
дополнительно подключать и настраивать отдельные расширения, например для
авторизации, работы с базой данных и миграциями. В Django значительная часть этих
механизмов уже предусмотрена, поэтому он лучше подходит для разрабатываемого
приложения.

Для разрабатываемого приложения особенно важна система аутентификации и
авторизации. Django содержит встроенную систему пользователей, групп, прав и
сессий \cite{django_auth}. Это позволяет реализовать роли, например
<<библиотекарь>> и <<читатель>>, а также ограничить доступ к отдельным действиям:
созданию и редактированию книг, оформлению выдач, просмотру статистики и
управлению читателями.

Для хранения данных выбрана SQLite. SQLite является самодостаточным SQL-движком,
не требующим отдельного серверного процесса и сложной настройки
\cite{sqlite_official}. База данных хранится в одном файле, что упрощает
развертывание и сопровождение приложения. Такой вариант подходит для небольшой
библиотеки, где количество пользователей ограничено, а операции записи выполняются
сравнительно редко.

Альтернативой SQLite является PostgreSQL. PostgreSQL представляет собой
полноценную объектно-реляционную систему управления базами данных с открытым
исходным кодом \cite{postgresql_official}. PostgreSQL лучше подходит для систем
с большим количеством пользователей, высокой нагрузкой, активной одновременной
записью данных и более сложным администрированием. Однако использование
PostgreSQL требует установки и сопровождения отдельного сервера баз данных,
создания пользователей, настройки прав доступа и резервного копирования. Для
небольшой библиотеки, которой требуется простое развертывание и базовый учет,
SQLite является более простым вариантом. При этом SQLite поддерживает основные
возможности реляционной базы данных, необходимые для приложения: создание таблиц,
хранение связанных данных, выполнение запросов, фильтрацию и агрегирование
данных.

Работа с базой данных будет выполняться через ORM Django. Это позволит описывать
модели приложения на языке Python, а затем автоматически создавать таблицы базы
данных с помощью миграций. При необходимости дальнейшего развития приложения
использование ORM упростит перенос проекта с SQLite на другую реляционную СУБД,
например PostgreSQL.

Для реализации пользовательского интерфейса используются HTML, CSS и JavaScript.
HTML задает структуру страниц, CSS отвечает за внешний вид, а JavaScript может
использоваться для небольших интерактивных элементов: подтверждения действий,
динамического отображения данных, улучшения форм и фильтров.

Для формирования HTML-страниц выбран подход с использованием серверных шаблонов
Django. Альтернативой мог бы быть отдельный frontend-фреймворк, например React
или Vue, при котором серверная часть предоставляет API, а пользовательский
интерфейс разрабатывается как отдельное клиентское приложение. Такой подход
целесообразен для более сложных интерактивных систем, но для приложения учета
книг, читателей и выдач он усложнил бы архитектуру. Использование шаблонов Django
позволяет реализовать интерфейс проще: сервер сразу формирует готовые HTML-страницы,
а JavaScript применяется только там, где требуется небольшая интерактивность.

Также в приложении предусматривается работа с файлами. Django содержит штатные
средства обработки файлов, загружаемых через формы, включая работу с объектом
\texttt{request.FILES} \cite{django_file_uploads}. В рамках приложения загрузка
файлов может использоваться для добавления обложек книг, а выгрузка данных ---
для формирования отчетов по книжному фонду и журналу выдач.

Сравнение выбранных инструментов и рассмотренных альтернатив приведено в
таблице~\ref{tab:dev_tools}.

{
\coursetablesetup
\setlength{\LTleft}{0pt}
\setlength{\LTright}{0pt}

\begin{longtable}{|
p{\dimexpr0.17\textwidth-2\tabcolsep-\arrayrulewidth\relax}|
p{\dimexpr0.17\textwidth-2\tabcolsep-\arrayrulewidth\relax}|
p{\dimexpr0.31\textwidth-2\tabcolsep-\arrayrulewidth\relax}|
p{\dimexpr0.35\textwidth-2\tabcolsep-\arrayrulewidth\relax}|}
    \caption{Сравнение инструментов разработки}
    \label{tab:dev_tools} \\

    \hline
    \textbf{Компонент} &
    \textbf{Выбранное решение} &
    \textbf{Рассмотренные альтернативы} &
    \textbf{Причина выбора} \\ \hline
    \endfirsthead

    \tabcont{\ref{tab:dev_tools}}
    \textbf{Компонент} &
    \textbf{Выбранное решение} &
    \textbf{Рассмотренные альтернативы} &
    \textbf{Причина выбора} \\ \hline
    \endhead

    \hline
    \endfoot

    Язык серверной части &
    Python &
    JavaScript на серверной стороне, C++ &
    Python удобен для реализации прикладной бизнес-логики и имеет развитую экосистему для веб-разработки \\ \hline

    Веб-фреймворк &
    Django &
    Flask &
    Django содержит встроенные механизмы пользователей, прав доступа, форм, ORM и административного интерфейса \\ \hline

    База данных &
    SQLite &
    PostgreSQL &
    SQLite проще в развертывании и сопровождении, так как не требует отдельного сервера БД и подходит для небольшой библиотеки \\ \hline

    Формирование интерфейса &
    HTML, CSS, шаблоны Django, JavaScript &
    Отдельный frontend-фреймворк, например React или Vue &
    Серверные шаблоны проще для приложения с типовыми страницами, формами и таблицами \\ \hline

    Работа с файлами &
    Средства Django для загрузки файлов &
    Ручная обработка файлов без встроенных механизмов фреймворка &
    Django предоставляет готовые средства обработки загруженных файлов через формы \\ \hline
\end{longtable}
}

Таким образом, выбранный стек технологий позволяет разработать веб-приложение
с понятной структурой, ролевой моделью доступа, интерфейсом управления данными,
поиском, фильтрацией, статистикой и обработкой файлов. При этом архитектура
остается достаточно простой для развертывания, сопровождения и возможного
дальнейшего расширения.

\subsection{Формулировка цели и задач работы}

По результатам анализа предметной области установлено, что для небольших библиотек
может быть востребовано веб-приложение, ориентированное не на комплексную
профессиональную автоматизацию всех библиотечных процессов, а на решение базовых
задач учета книг, читателей и выдач. Такое приложение должно быть проще во
внедрении и сопровождении, чем полнофункциональные библиотечно-информационные
системы, но при этом должно закрывать основные операции ежедневной работы
библиотеки.

Целью разработки является создание веб-приложения для автоматизации учета книг
и читателей библиотеки, обеспечивающего ведение книжного фонда, учет читателей,
регистрацию выдачи и возврата книг, поиск информации и формирование статистики.

Для достижения указанной цели необходимо решить следующие задачи:

\begin{enumerate}
    \item определить основные процессы предметной области и состав пользователей
    приложения;

    \item спроектировать модель данных для хранения информации о книгах,
    читателях, авторах, категориях и выдачах;

    \item реализовать механизм регистрации, входа в систему и разграничения
    доступа пользователей по ролям;

    \item разработать интерфейсы для добавления, просмотра, редактирования и
    удаления записей о книгах, читателях и выдачах;

    \item реализовать поиск и фильтрацию книг по основным параметрам;

    \item реализовать обработку данных и формирование статистики по книжному
    фонду и журналу выдач;

    \item реализовать загрузку и выгрузку файлов, связанных с данными приложения;

    \item выполнить проверку работоспособности основных функций веб-приложения.
\end{enumerate}

В качестве основных сущностей приложения предполагается использовать следующие:
\texttt{Book} --- книга, \texttt{Reader} --- читатель, \texttt{Loan} --- запись
о выдаче книги. Дополнительно могут быть выделены сущности \texttt{Author},
\texttt{Category} и \texttt{UploadedFile}. Такая структура данных позволяет
описать основные объекты предметной области и связать учет книг с действиями
читателей.